% !TeX root = ../main.tex
\section{泥石流触发临界条件与风险评估（问题二）}\label{sec:thresholds}

\subsection{冰岩崩—泥石流触发机制分析}\label{sec:trigger-mechanisms}
% 区分冰体初始失稳和泥石流形成机制，建立静态条件与动态触发因素的关系。
\writingplaceholder

\subsection{基于地形与物源的空间易发性评估}\label{sec:susceptibility}
% 利用坡度、地形起伏、冰川几何、物源和土地覆盖等特征建立空间易发性模型。
\writingplaceholder

\subsection{基于气象与冰川状态的动态触发指标}\label{sec:dynamic-indicators}
% 建立随时间变化的气象、水热与冰体状态指标，考虑不同物理过程的作用。
\writingplaceholder

\subsection{关键参数临界条件识别}\label{sec:critical-conditions}
% 研究单因素阈值及多因素联合临界边界，区分统计决策阈值与具有物理依据的失稳条件。
\writingplaceholder

\subsection{时空耦合触发风险评估模型}\label{sec:coupled-risk}
% 将静态易发性和动态触发条件结合，输出不同子流域、不同时间窗口的风险评价。
\writingplaceholder

\subsection{模型时间与空间可靠性分析}\label{sec:risk-reliability}
% 比较不同时间尺度、不同空间分辨率、不同子流域的结果；在具备事件标签的情况下进行独立验证和校准。
\writingplaceholder
